\documentclass[11pt,a4paper]{article}

\usepackage{amsmath}
\usepackage{amssymb}
\usepackage{amsthm}
\usepackage{geometry}
\usepackage{xcolor}
\usepackage{booktabs}
\usepackage[most]{tcolorbox}
\usepackage{hyperref}

\geometry{margin=1in}

\hypersetup{
    colorlinks=true,
    linkcolor=blue,
    filecolor=magenta,      
    urlcolor=cyan,
    citecolor=red,
}

% Theorem environments for strict SCT structure
\newtheorem{theorem}{Theorem}[section]
\newtheorem{proposition}[theorem]{Proposition}
\newtheorem{lemma}[theorem]{Lemma}
\newtheorem{corollary}[theorem]{Corollary}

\theoremstyle{definition}
\newtheorem{definition}[theorem]{Definition}
\newtheorem{convention}[theorem]{Convention}

\theoremstyle{remark}
\newtheorem{remark}[theorem]{Remark}

% Epistemic status tags
\newcommand{\exact}{\textcolor{blue}{[exact from lit.] }}
\newcommand{\verified}{\textcolor{green!60!black}{[verified] }}
\newcommand{\sctpred}{\textcolor{red}{[SCT prediction] }}
\newcommand{\gap}{\textbf{\textcolor{red}{[GAP]}} }

\begin{document}

\title{\textbf{NT-Test: SCT Form Factors and Propagator Criteria} \\ \large Formatting and Style Verification}
\author{David Alfyorov}
\date{March 2026 --- Working Document}

\maketitle


\begin{center}
\small\textit{Credits: Aliaksandr Samatyia contributed research-assistance and workflow support.}
\end{center}

\begin{tcolorbox}[colback=green!5!white,colframe=green!60!black,title=\textbf{STATUS: VERIFIED}]
This document serves as the new formatting baseline for SCT Theory derivations, incorporating the six-stage review layout, explicit epistemic tags, and theorem environments.
\end{tcolorbox}

\begin{abstract}
We establish the standard formatting for SCT Theory derivations by summarizing the entire-function structure of the nonlocal form factors $F_1(z)$ and $F_2(z, \xi)$ together with the propagator criteria that must be checked separately in the physical denominator analysis.
\end{abstract}

\tableofcontents
\vspace{1cm}

\section{Setup and Conventions}

\begin{convention}[Metric signature]
\exact We use the mostly-plus Lorentzian signature $\eta_{\mu\nu} = \text{diag}(-1, +1, +1, +1)$. For heat kernel computations, we Wick-rotate to Euclidean signature.
\end{convention}

\begin{convention}[Dirac Operator Data]
\verified For the massless Dirac operator on a Riemannian 4-manifold, the generalized Laplacian $\Delta = -D^2 - E$ has the following data:
\begin{equation}
    \text{tr}(\mathbf{1}) = 4\,, \quad E = -\frac{R}{4}\,, \quad \Omega_{\mu\nu} = \frac{1}{4}R_{\mu\nu\rho\sigma}\gamma^\rho\gamma^\sigma\,.
\end{equation}
\end{convention}

\section{Analyticity and Pole Cancellation}

\begin{definition}[Master function]
The nonlocal form factors are built from the master function:
\begin{equation}
    \phi(z) = \int_0^1 \mathrm{d}\alpha\, e^{-\alpha(1-\alpha)z}\,.
\end{equation}
\end{definition}

\begin{proposition}[Taylor expansion]
\verified The master function has the expansion $\phi(z) = 1 - \frac{z}{6} + \frac{z^2}{60} + \mathcal{O}(z^3)$, which yields $\phi(0)=1$ and $\phi'(0)=-1/6$.
\end{proposition}

\begin{theorem}[Entire-function status of $F_1$]
\sctpred The total Standard Model Weyl form factor $F_1(z)$ is an entire function.
\end{theorem}

\begin{proof}
The apparent poles in the Codello--Zanusso form factors are of the form $1/z$ and $(\phi-1)/z^2$. Using the Taylor expansion from the previous proposition, the numerator combinations identically vanish at $z=0$, rendering the singularity removable. 
\end{proof}

\section{Results and Benchmarks}

\begin{tcolorbox}[colback=blue!5!white,colframe=blue!75!black,title=\textbf{Key Result: SM Local Limits}]
The local limits of the total SM form factors are:
\begin{align}
    \alpha_C &= 16\pi^2 F_1(0) = \frac{13}{120}\,, \\
    \alpha_R(\xi) &= 16\pi^2 F_2(0, \xi) = 2\left(\xi - \frac{1}{6}\right)^2\,.
\end{align}
\end{tcolorbox}

\begin{table}[h]
\centering
\begin{tabular}{lllc}
\toprule
\textbf{Quantity} & \textbf{Value} & \textbf{Source} & \textbf{Status} \\
\midrule
$h_C^{(0)}(0)$ & $1/120$ & CZ Eq. (18) & \verified \\
$h_C^{(1/2)}(0)$ & $-1/20$ & CZ Eq. (16) & \verified \\
$\alpha_C$ & $13/120$ & NT-1b Phase 3 & \verified \\
\bottomrule
\end{tabular}
\caption{Benchmark local limits for form factors.}
\end{table}

\end{document}
